% 01_introduction -- written by Claude Code (D-045); revised after the peer review (Prompt 16: M1, M3, M5, m19).
\section{Introduction}\label{sec:intro}
AI systems, from rule-based decision support to large language models (LLMs), are increasingly evaluated against the
decisions of multidisciplinary tumor boards~\cite{sharma2026review,zhao2020watson,liao2025,li2025imitation,asoglu2026},
and several recent studies describe a single model recommendation per case~\cite{popescu2026,zhou2026lung,dindelegan2026}.
Yet an LLM asked the same question again can answer differently: where tumor-board studies repeated their queries,
agreement ranged from about a third of cases for ChatGPT~\cite{liao2025,asoglu2026} to near-perfect for some models and
poor for reasoning models~\cite{li2025imitation}, and on benchmarks consistency also varies with the
temperature~\cite{pinhanez2025,liu2026repro}. Repetition can also be used: several sampled answers can be combined by
a majority vote (self-consistency~\cite{wang2022selfconsistency}), and agreement between runs can signal how far an
answer can be trusted~\cite{rydzewski2024}. In oncology, voting has been tested with strong closed models and brought
little benefit~\cite{holmes2023,stifini2026,windisch2026}; whether it helps a small open model, which varies more,
is not known.

We use oncology questions from the MedQA licensing-exam dataset~\cite{jin2021medqa} as a controlled testbed with a
single keyed answer, and compare greedy decoding with a majority vote over 10 sampled runs of Qwen2.5-7B-Instruct on
\nTest{} questions, under a protocol frozen before any model output on test items existed. We ask: (RQ1) Is the
majority of 10 sampled runs more accurate than the greedy answer? (RQ2) How consistent are the answers across runs?
(RQ3) Does agreement between runs flag answers that are more likely to be correct? Our contributions are:
\begin{itemize}
\item a pre-specified paired comparison in which the vote was \diffMain{} points more accurate than greedy decoding
  (95\% CI \diffMainLo{} to \diffMainHi) at \tokRatio{} times the output tokens, \decompVote{} points of it, post hoc,
  attributable to voting itself;
\item like-for-like consistency figures: greedy repeated its answer for \rerunSamePct\% of items, two sampled runs
  agreed for \pairAgree\%, all ten for \unanimity\%;
\item evidence that agreement tracks correctness (AUROC \auroc), with its limits: it partly reflects item difficulty,
  \unanimousWrongPct\% of unanimous answers were wrong, and the runs share an answer-position preference;
\item a frozen protocol, code, item identifiers and parsed answers from which every statistic can be recomputed.
\end{itemize}
